% Complete editorial continuation arguments; frozen preceding packet is unchanged.
\section{Invariance, recursive universal properties and relations}

\begin{annexitem}{annex:p10.domain-invariance.least-hom}{thm:p10-domain-invariance-least-hom}{Least comparisons exist and compose}
\begin{proof}
Local continuity of \(F\) and continuity of composition make
\(\phi_{C,A}\) continuous on the pointed function dcpo.
Its fixed-point equation is equivalent to \(ui=aF(u)\), since
\(i\) is invertible. Thus its least fixed point is exactly the
least homomorphism, by the proved least-fixed-point theorem.
For the composition assertion define
\(H:[B\longrightarrow C]\to[A\longrightarrow C]\) by
\(H(v)=vh_{B,A}\). This is continuous and strict as a map of
function dcpo's: precomposing a constant-bottom map remains
constant bottom even when \(h_{B,A}\) itself is non-strict.
The homomorphism equation for \(h_{B,A}\) gives
\(j^{-1}h_{B,A}=F(h_{B,A})i^{-1}\), so
\[
 H(\phi_{C,B}(v))=aF(v)j^{-1}h_{B,A}
     =aF(vh_{B,A})i^{-1}=\phi_{C,A}(H(v)).
\]
Strict uniformity of least fixed points gives the asserted
composition equality. On a solution itself identity is a fixed
point of \(\phi_{A,A}\), hence \(\mu_A\sqsubseteq1_A\).
Applying the composition equality with all three solutions equal
gives \(\mu_A^2=\mu_A\). Since \(\mu_A\) is below identity,
its value at bottom must be bottom, proving strictness.
\end{proof}
\end{annexitem}

\begin{annexitem}{annex:p10.domain-invariance.covariant}{thm:p10-domain-invariance-covariant}{Canonical approximants establish the identity and the comparison pair}
\begin{proof}
The seed embedding \(e_0:I\to D\) selects bottom, so
\(q_0=e_0p_0=c_{\bot}\). The two fold leg equations in the
frozen construction give
\[
 \theta_D(q_n)=i_DF(e_n)F(p_n)i_D^{-1}
             =e_{n+1}p_{n+1}=q_{n+1}.
\]
Consequently its iterates from bottom are the actual bilimit
approximants. Their supremum is identity, so \(\mu_D=1_D\).
The least-homomorphism composition theorem gives
\(pe=h_{D,A}h_{A,D}=\mu_D=1_D\) and
\(ep=\mu_A\sqsubseteq1_A\); this proves minimality with
actual maps and the correct information-order inequality.

If \(\mu_A=1_A\), these comparisons are inverse algebra
homomorphisms, hence a solution isomorphism. Conversely a
solution isomorphism \(t:D\to A\) conjugates endomaps by
\(u\mapsto tut^{-1}\). It is strict and continuous on the
function dcpo and intertwines \(\theta_D,\theta_A\), by its
structure equation and functoriality. Uniformity carries their
least fixed points to one another, giving \(\mu_A=1_A\).
For a strict algebra endomorphism \(h\), postcomposition
\(H(u)=hu\) is a strict continuous function-dcpo map and
\(H\theta_A=\theta_AH\) follows from \(hi=iF(h)\).
Uniformity gives \(H(\mu_A)=\mu_A\). If \(\mu_A=1_A\),
this says \(h=1_A\). Conversely \(\mu_A\) is itself a
strict endomorphism, so uniqueness forces it to be identity.
Every step also applies in the strict-only category: the
approximants, isomorphisms and hom maps used there are strict,
and the strict function dcpo has the same pointwise chain suprema.
No unrestricted extension of a strict-only functor is needed.
\end{proof}
\end{annexitem}

\begin{annexitem}{annex:p10.domain-invariance.initial-final}{thm:p10-domain-invariance-initial-final}{Why initiality needs strictness and finality has a stronger scope}
\begin{proof}
First take the strict category. For a strict structure
\(a:F(A)\to A\), the functional
\(\phi(u)=aF(u)\operatorname{unfold}\) maps strict homs to
strict homs and is continuous on that pointed hom dcpo.
Its least fixed point exists and is a strict algebra homomorphism.
For any other strict homomorphism \(h:D\to A\), let
\(H:[D\mathrel{\mathop{\longrightarrow}^{\bot!}}D]\to
[D\mathrel{\mathop{\longrightarrow}^{\bot!}}A]\) be \(H(v)=hv\).
It is continuous and strict because \(h\) is strict. Its
homomorphism equation gives \(H\theta_D=\phi H\).
Since \(\operatorname{fix}(\theta_D)=1_D\), strict uniformity
gives \(h=H(1_D)=\operatorname{fix}(\phi)\).
This proves existence and uniqueness in the asserted category.

Now take an unrestricted locally continuous \(F\) and an
arbitrary continuous coalgebra \(c:A\to F(A)\). The functional
\(\psi(u)=\operatorname{fold}F(u)c\) is continuous on
\([A\longrightarrow D]\), so has a least fixed point.
Multiplying its equation by unfold shows that this is a coalgebra
homomorphism. For any such homomorphism \(h\), use
\(H:[D\longrightarrow D]\to[A\longrightarrow D]\),
\(H(v)=vh\). This map is strict on function dcpo's irrespective
of whether \(h\) is strict, since the constant-bottom function
remains constant bottom after precomposition. The coalgebra
equation \(\operatorname{unfold}h=F(h)c\) gives
\[
 H\theta_D(v)=\operatorname{fold}F(v)\operatorname{unfold}h
            =\operatorname{fold}F(vh)c=\psi H(v).
\]
Uniformity again gives \(h=H(1_D)=\operatorname{fix}(\psi)\).
This is the promised uniqueness for non-strict coalgebras too.
For a functor defined only on strict maps, take \(c\) strict
and work on strict hom dcpo's; the same proof applies, but does
not assert that \(F(h)\) is defined for a non-strict \(h\).
\end{proof}
\end{annexitem}

\begin{annexitem}{annex:p10.domain-invariance.least-pairs}{thm:p10-domain-invariance-least-pairs}{The composition assertion includes every approximation stage}
\begin{proof}
Joint local continuity and composition make \(\Phi_{B,A}\)
continuous on the pointed product function dcpo. Its fixed points
are precisely the two displayed comparison equations, after
multiplication by \(i,j\). The least-fixed-point theorem therefore
gives the least comparison pair.
Write \((h^n_{B,A},k^n_{A,B})\) for its iterates from the two
constant-bottom maps. For solutions \(A,B,C\), induction gives
\[
 h^n_{C,B}h^n_{B,A}=h^n_{C,A},\qquad
 k^n_{A,B}k^n_{B,C}=k^n_{A,C}.
\]
At zero, the two composites are the appropriate constant-bottom
maps. At the next stage, the forward composite is
\[
 \ell F(k^n_{B,C},h^n_{C,B})j^{-1}
       jF(k^n_{A,B},h^n_{B,A})i^{-1}
 =\ell F(k^n_{A,B}k^n_{B,C},
                 h^n_{C,B}h^n_{B,A})i^{-1}.
\]
It is the next forward approximation for \(A,C\); the backward
calculation is its correctly reversed counterpart. Passing to
suprema is justified by joint continuity of composition. The
two independent approximation indices have cofinal diagonal,
so the equalities pass to the least comparison pairs themselves.
This proves the corrected-author composition comment explicitly.
With \(A=B\), the two iterates are equal by symmetry and follow
\(u\mapsto iF(u,u)i^{-1}\); hence both limits equal \(\mu_A\).
The pair \((1_A,1_A)\) is fixed, so \(\mu_A\sqsubseteq1_A\),
and the composition equality gives \(\mu_A^2=\mu_A\).
Being below identity forces its bottom value to bottom.
\end{proof}
\end{annexitem}

\begin{annexitem}{annex:p10.domain-invariance.mixed}{thm:p10-domain-invariance-mixed}{Mixed approximants, minimality and the extra order step in uniqueness}
\begin{proof}
For the canonical mixed solution, \(q_0=c_{\bot}\). Its leg
maps are \(F(p_n,e_n)\) with projections \(F(e_n,p_n)\).
Therefore
\[
 \theta_D(q_n)=\operatorname{fold}F(e_np_n,e_np_n)
                     \operatorname{unfold}
 =\operatorname{fold}F(p_n,e_n)F(e_n,p_n)
                     \operatorname{unfold}=q_{n+1}.
\]
The canonical identity-supremum equation gives \(\mu_D=1_D\).
Composing the two least comparison pairs gives identity on \(D\)
and \(\mu_A\) on \(A\), so they are an e-p-pair. If
\(\mu_A=1_A\), they are inverse solution comparisons.
Conversely a solution isomorphism \(t:D\to A\), with reverse
map \(t^{-1}\), conjugates the two mixed endomap functionals:
the equations \(ti_D=i_AF(t^{-1},t)\) and their inverses,
with contravariance, give
\(t\theta_D(u)t^{-1}=\theta_A(tut^{-1})\).
This conjugation is strict and continuous on function dcpo's;
uniformity thus carries \(\mu_D=1_D\) to \(\mu_A=1_A\).

Suppose now \(\mu_A=1_A\) and \(h\) is a strict mixed
endomorphism. Let \(H(u)=huh\). It is continuous and strict
on the endomap dcpo. Using
\(hi=iF(h,h)\) and \(i^{-1}h=F(h,h)i^{-1}\),
\[
 H\theta_A(u)=iF(h,h)F(u,u)F(h,h)i^{-1}
             =iF(huh,huh)i^{-1}=\theta_AH(u).
\]
Uniformity gives \(h^2=H(1)=1\). This alone is not the
desired conclusion. Since \(h\) is a fixed point of
\(\theta_A\) and its least fixed point is identity, also
\(1\sqsubseteq h\). Monotonicity of \(h\) applied to this
pointwise inequality gives \(h\sqsubseteq h^2=1\), so
\(h=1\). Conversely \(\mu_A\) is a strict mixed
endomorphism; uniqueness therefore forces \(\mu_A=1\).
The calculation on strict-only hom dcpo's uses only strict maps
and proves the same canonical identity in that regime as well.
\end{proof}
\end{annexitem}

\begin{annexitem}{annex:p10.domain-invariance.free-mixed}{thm:p10-domain-invariance-free-mixed}{Every solution pair is forced by the canonical approximants}
\begin{proof}
The stated functional \(\Psi\) is well typed on the strict
product hom dcpo and continuous by local continuity and
composition. Strictness of \(f,g\), fold, unfold and all
maps in its functor regime makes every output pair strict.
Its least fixed point therefore exists and satisfies the two
universal equations, after composing with the inverse structure maps.
Let \((h,k)\) be any strict pair satisfying those equations.
The canonical mixed identity calculation still holds here, since
it uses only the strict approximants \(q_n\):
\(q_0=c_{\bot}\), \(q_{n+1}=\operatorname{fold}F(q_n,q_n)
\operatorname{unfold}\), and \(\bigsqcup_n^{\uparrow}q_n=1\).
We claim that the \(n\)-th iterate of \(\Psi\) from bottom is
\((q_nh,kq_n)\). At zero this is the bottom pair; the backward
component uses the actual strictness \(k(\bot_D)=\bot_B\).
For the next forward component, functoriality gives
\[
 \operatorname{fold}F(kq_n,q_nh)f
 =\operatorname{fold}F(q_n,q_n)F(k,h)f
 =q_{n+1}h.
\]
For the backward component it gives
\[
 gF(q_nh,kq_n)\operatorname{unfold}
 =gF(h,k)F(q_n,q_n)\operatorname{unfold}=kq_{n+1}.
\]
Continuity of composition now makes the supremum of these iterates
\((h,k)\). Thus every solution pair equals the same least fixed
point of \(\Psi\), proving uniqueness rather than only leastness.
\end{proof}
\end{annexitem}

\begin{annexitem}{annex:p10.domain-invariance.parameters}{thm:p10-domain-invariance-parameters}{Parameter maps commute stage by stage and preserve directed suprema}
\begin{proof}
Let \(a_n^A:X_n(A)\to X_{n+1}(A)\) be the adjacent
canonical embeddings. The maps \(u_n\) are strict. They satisfy
\(u_{n+1}a_n^A=a_n^{A'}u_n\): at zero both sides are
the unique strict map out of \(I\); at the next stage apply
functoriality of \(F\) to the preceding equation and the
parameter identities. Hence \(e_n^{A'}u_n\) is an actual
cocone over the source sequence.
For \(n\le m\), its stage maps obey
\(u_ma_{mn}^A=a_{mn}^{A'}u_n\). Thus
\[
 e_n^{A'}u_np_n^A
 =e_m^{A'}u_ma_{mn}^Ap_{nm}^Ap_m^A
                   \sqsubseteq e_m^{A'}u_mp_m^A.
\]
Their directed supremum is the bilimit cocone mediator, and is
strict because each displayed map is strict. Its leg equation is
\(M(u)e_n^A=e_n^{A'}u_n\).
For identity parameters all stage maps are identity; for a
composite \(vu\), induction and functoriality give
\((vu)_n=v_nu_n\). Identity and composition of the induced
maps follow from their leg equations and colimit uniqueness.

For a directed parameter family \(U\), induction gives
\((\bigsqcup^{\uparrow}U)_n=\bigsqcup_{u\in U}^{\uparrow}u_n\).
At the successor stage local continuity initially supplies two
indices, one for the parameter map and one for its earlier-stage
map. Their diagonal is cofinal, since the latter assignment is
monotone by the same induction. This gives precisely the single-
index formula, not an assumption that a dependent index can be
discarded. Continuity of composition and interchange of the two
directed suprema over \(n,U\) proves local continuity of \(M\).

To prove naturality, compose its proposed fold equation with
every colimit leg \(F(1_A,e_n^A)\). Its left route is
\(M(u)e_{n+1}^A=e_{n+1}^{A'}u_{n+1}\).
Its right route is
\(\alpha_{A'}F(u,e_n^{A'}u_n)
=e_{n+1}^{A'}F(u,u_n)\), the same map. Since these are
colimit legs, they jointly determine maps out of their colimit,
proving the equation. Multiplying by the two inverse folds gives
the stated unfold equation; no nonexistent adjoint to \(M(u)\)
is invoked.
\end{proof}
\end{annexitem}

\begin{annexitem}{annex:p10.domain-invariance.mixed-parameters}{thm:p10-domain-invariance-mixed-parameters}{Four typed stage compatibilities make the parameter functor real}
\begin{proof}
For fixed \(P,Q\), the two recursive coordinates are the
diagonal of the mixed functor on pairs
\[
 ((U,V),(X,Y))\longmapsto
       (F(P,Q,V,X),\ F(Q,P,U,Y)).
\]
Its hom action is jointly continuous, coordinate by coordinate.
The embedding passage of the frozen packet therefore supplies
actual e-p-pairs at every stage and preserves the coordinatewise
bilimit. Swapping \(P,Q\) swaps the two whole stage diagrams;
their actual tuple limits give the indicated \(M(P,Q),M(Q,P)\)
and the two fold isomorphisms.

Write \(a_n^X,b_n^X\), \(a_n^Y,b_n^Y\) for adjacent
embeddings/projections, and use primes for the \(P',Q'\)
diagram. The displayed parameter stage maps satisfy all four
identities
\[
 \begin{aligned}
 r_{n+1}a_n^X&=a_n^{X'}r_n,& b_n^{X'}r_{n+1}&=r_nb_n^X,\\
 s_{n+1}a_n^{Y'}&=a_n^Ys_n,& b_n^Ys_{n+1}&=s_nb_n^{Y'}.
 \end{aligned}
\]
At zero, the embedding equations are equality of strict maps
out of \(I\), and the projection equations are equality of
maps into \(I\). At a later stage the first equation follows
by expanding its left side as
\[
 F(u,v,s_n,r_n)F(1,1,b_{n-1}^Y,a_{n-1}^X)
 =F(u,v,b_{n-1}^Ys_n,r_na_{n-1}^X),
\]
then using the previous projection compatibility for \(s\)
and embedding compatibility for \(r\). It becomes
\(F(u,v,s_{n-1}b_{n-1}^{Y'},a_{n-1}^{X'}r_{n-1})\),
the expansion of the right side. For the second equation the
left side expands as
\[
 F(1,1,a_{n-1}^{Y'},b_{n-1}^{X'})F(u,v,s_n,r_n)
 =F(u,v,s_na_{n-1}^{Y'},b_{n-1}^{X'}r_n),
\]
which the other two preceding compatibilities turn into
\(F(u,v,a_{n-1}^Ys_{n-1},r_{n-1}b_{n-1}^X)\),
the right side. The same two calculations, with source/target
exchanged and \(u,v\) swapped, are exactly the third and
fourth equations, of types \(Y'_n\to Y_n\). Thus all
four follow by simultaneous induction, not by assuming parameter
maps are embedding--projection pairs.

The embedding compatibilities give the two cocones and the
increasing mediator families as in the positive-parameter proof.
The projection compatibilities additionally give
\(p_n^{X'}r=r_np_n^X\) and
\(p_n^Ys=s_np_n^{Y'}\): evaluate the defining supremum on
the cofinal tail of stages \(m\ge n\), where each expression
is constant by those compatibilities.
For composable parameter pairs, induction and functoriality
give composition of the \(r_n\)'s in the forward order and
of the \(s_n\)'s in the reversed order. Identities are stage
identities. Colimit uniqueness gives the mixed functor laws.
For directed parameter-map families, simultaneous induction
through the two recurrences shows that both stage assignments
preserve directed suprema; local continuity of \(F\) and a
common upper pair of indices give the diagonal step. Continuity
of composition and interchange with the stage supremum then
give local continuity of the two limit actions.

Finally compose the first naturality equation with the fold
legs \(F(1,1,p_n^Y,e_n^X)\). On the left it gives
\(e_{n+1}^{X'}r_{n+1}\). On the right use
\(p_n^Ys=s_np_n^{Y'}\) and \(re_n^X=e_n^{X'}r_n\),
then the target fold leg, obtaining the same map. The legs
are colimiting, hence determine the equation. The second
coordinate is the exchanged equation. Inverting the folds
gives naturality of the two unfolds. All negative slots have
actual backward maps; no negative slot was filled with a
forward map of the wrong type.
\end{proof}
\end{annexitem}

\begin{annexitem}{annex:p10.domain-invariance.bekic}{thm:p10-domain-invariance-bekic}{Least pairs and the nested initial algebra have the same universal meaning}
\begin{proof}
The continuous fixed-point operator makes \(h\) continuous,
and pairing/composition make \(k\) continuous. The pair
\((h(d),d)\) satisfies both fixed-point equations. For any
joint pre-fixed point \((x,y)\), leastness gives \(h(y)\sqsubseteq x\).
Thus \(k(y)=g(h(y),y)\sqsubseteq g(x,y)\sqsubseteq y\),
so \(d\sqsubseteq y\) and
\(h(d)\sqsubseteq h(y)\sqsubseteq x\). The nested pair
is therefore the least joint pre-fixed point, and equals the
simultaneous least fixed point. This argument includes non-strict
continuous maps and is not a categorical initiality claim about them.

For the domain assertion, the parameter theorem supplies the
locally continuous strict \(H\) and its natural folds
\(\alpha_B:P(H(B),B)\to H(B)\).
Then \(K(B)=Q(H(B),B)\) is locally continuous, so its
canonical fold \(\beta:K(B_*)\to B_*\) is initial.
Given any strict simultaneous algebra
\(a:P(A,B)\to A,b:Q(A,B)\to B\), initiality of
\(\alpha_B\) supplies the unique
\(u_B:H(B)\to A\) with
\(u_B\alpha_B=aP(u_B,1_B)\).
Make \(B\) a \(K\)-algebra by
\(t=bQ(u_B,1_B)\). Initiality of \(\beta\) supplies
unique \(v:B_*\to B\) with \(v\beta=tK(v)\).
Set \(u=u_BH(v)\). Naturality of \(\alpha\) gives
\[
 u\alpha_{B_*}
 =u_B\alpha_BP(H(v),v)=aP(u_BH(v),v)=aP(u,v),
\]
and the equation for \(v\) is \(v\beta=bQ(u,v)\).
Thus \((u,v)\) is an actual simultaneous algebra morphism.

If \((h,k)\) is another one, view \(A\) as a
\(P(-,B_*)\)-algebra through \(aP(1_A,k)\).
Both \(h\) and \(u_BH(k)\) are homomorphisms from
\(\alpha_{B_*}\), the latter by naturality; initiality gives
\(h=u_BH(k)\). Its other equation becomes
\(k\beta=tK(k)\), so initiality of \(\beta\) gives
\(k=v\), then \(h=u\). This proves nested initiality.
The simultaneous canonical construction in the product category
is also initial by the already proved strict initiality theorem:
its hom dcpo's, seed and bilimit are coordinatewise. The unique
homomorphisms between the two initial algebras are inverse,
since each composite is an endomorphism of an initial algebra
and must be identity. This supplies the stated structure
isomorphism rather than only matching object names.
\end{proof}
\end{annexitem}

\begin{annexitem}{annex:p10.domain-invariance.resumptions}{thm:p10-domain-invariance-resumptions}{Strict constructor action and an infinite observable resumption}
\begin{proof}
For strict \(u:R\to R'\), the map \(1_S\times u\) is
strict and continuous, so the coalesced-sum universal property
gives the actual strict continuous map
\(J_u=1_S\oplus(1_S\times u)\).
Postcomposition gives a continuous map \(T(u)\) on the
continuous function dcpo. It is strict there because \(J_u\)
sends the sum's bottom to bottom. Identity and composition
follow first in the two sum branches and then by postcomposition.
For a directed family \(U\) of strict maps, on a final-state
branch \(J_u\) is independent of \(u\); on a continuation
branch its value is \(\operatorname{inr}(s,u(r))\).
Continuous pairing and the continuous branch injection give
\(J_{\bigsqcup U}=\bigsqcup_{u\in U}J_u\) pointwise,
including when the product bottom is collapsed.
Consequently \(T(\bigsqcup U)=\bigsqcup_{u\in U}T(u)\)
pointwise. This verifies local continuity in its correct strict
regime. The input computations \(r:S\to S\oplus(S\times R)\)
are not required strict in their state input; strictness here
concerns maps between the resumption objects.
Canonical mixed/ordinary embedding preservation in this positive
case yields the fold/unfold construction. The strict initial
and strict final theorems apply to this actual functor. No
statement here requires \(T(u)\) for a non-strict \(u\).

The map \(r\mapsto(x\mapsto\operatorname{inr}(s,r))\)
is continuous by pairing, branch injection and pointwise function
suprema. Composing with fold gives the continuous step functional.
Its iterates \(r_0=\bot\), \(r_{n+1}=\operatorname{step}(s,r_n)\)
are increasing. Unfolding \(r_n\) for \(n>0\) gives the
constant continuation \((s,r_{n-1})\); induction on the
number of observations exposes exactly \(n\) successive
states before the unknown remainder \(r_0\).
Let \(\Omega_s=\bigsqcup_n^{\uparrow}r_n\).
Continuity of unfold, pairing and injection gives
\[
 \operatorname{unfold}(\Omega_s)(x)
                  =\operatorname{inr}(s,\Omega_s)
\]
for every input state \(x\). Since \(s\ne\bot\), this
is a nonbottom continuation, never a final branch.
In contrast \(\operatorname{unfold}(\operatorname{return}_s)(x)
=\operatorname{inl}(s)\). The two nonbottom tags are different,
so these resumptions are distinct. These are actual semantic
observations, not claims about termination of an unspecified
operational language.
\end{proof}
\end{annexitem}

\begin{annexitem}{annex:p10.domain-invariance.relational}{thm:p10-domain-invariance-relational}{An admissible predicate on maps carries the general comparison proof}
\begin{proof}
The constant-bottom map sends every tuple in \(R\) to the
constantly-bottom tuple in \(S\), so belongs to the map predicate.
If \(u_m\) is an increasing omega-chain in that predicate,
then for each \(\bar x\in R\), the tuples \(u_m\bar x\)
form a chain in \(S\). Admissibility puts their coordinatewise
supremum \((\bigsqcup_m u_m)\bar x\) in \(S\).
The supremum map is strict and continuous by the earlier hom
dcpo theorem, so the predicate is admissible there.

The source action condition implies its variance assertions by
taking identity maps: a negative inclusion \(R_2\subseteq R_1\)
gives \(F^{\mathrm{rel}}(R_1,S)\subseteq
F^{\mathrm{rel}}(R_2,S)\); a positive inclusion gives the
same direction. For canonical \(D\) let
\(P=\{u:u^{\mathrm{rel}}R\subseteq S\}\).
For \(u\in P\), the two assumed inclusions and the typed
action compatibility give
\[
 \begin{aligned}
 \theta_D(u)^{\mathrm{rel}}R
 &=\operatorname{fold}^{\mathrm{rel}}F(u,u)^{\mathrm{rel}}
                 \operatorname{unfold}^{\mathrm{rel}}R\\
 &\subseteq\operatorname{fold}^{\mathrm{rel}}F(u,u)^{\mathrm{rel}}
                       F^{\mathrm{rel}}(S,R)\\
 &\subseteq\operatorname{fold}^{\mathrm{rel}}F^{\mathrm{rel}}(R,S)
 \subseteq S.
 \end{aligned}
\]
In the middle application, the negative map \(u\) takes \(R\)
into \(S\), and the positive map \(u\) takes \(R\) into
\(S\), exactly the two inputs required by the action definition.
Thus \(P\) is invariant under \(\theta_D\). Fixed-point
induction puts its least fixed point, identity, in \(P\).
This is \(R\subseteq S\). For arity zero the relation
containing the bottom empty tuple is already the whole singleton,
and the same argument or direct inclusion applies.
\end{proof}
\end{annexitem}

\begin{annexitem}{annex:p10.domain-invariance.actions}{thm:p10-domain-invariance-actions}{Every primitive action has the claimed carrier, admissibility and compatibility}
\begin{proof}
For unary actions, lifting contains new bottom; an increasing
chain is either entirely bottom or has a cofinal old part whose
values lie in \(R\), so its supremum remains in the lifted
predicate. Products use coordinatewise chain suprema. For sums,
the nonbottom part of an increasing chain lies in one branch;
its supremum remains there by the component predicate's closure.
For smashed products use the same old-pair argument with both
nonbottom coordinates, or the bottom-only case. For a chain of
functions mapping \(R\) into \(S\), pointwise closure of \(S\)
shows that its supremum does too. Its bottom function has value
bottom in \(S\). Restricting to strict functions preserves this
argument. These observations prove admissibility of every unary
primitive action.

For binary lifting the exact relation is
\[
 \{(x,y)\in(D_{\bot})^2:x=\bot_{D_{\bot}}\}\quad\cup\quad
 \{(\operatorname{up}a,\operatorname{up}b):(a,b)\in R\}.
\]
It contains the bottom pair. If the supremum of an increasing
chain has old first coordinate, a cofinal tail has both coordinates
old and related by \(R\); admissibility of \(R\) handles its
supremum. If its first coordinate remains new bottom, it is
already in the first clause. Products, sums and smash use the
same coordinate/branch classification; when the first value has
become nonbottom, all later related second values lie in the
same indicated branch or old product, and component closure
applies. For binary function actions, at each fixed argument
the two function chains give a pair chain in \(S\), so pointwise
suprema remain related. The bottom pair of functions is included.

Compatibility for lifting and tagged/paired constructors follows
by applying the given relation-preserving component maps; strict
maps also preserve the bottom clause and the collapsed axes.
For unary function space, if \(f(R_1)\subseteq S_1\),
\(u(R_2)\subseteq R_1\), \(v(S_1)\subseteq S_2\), then
\(vfu(R_2)\subseteq S_2\), giving the correct contravariant
and covariant condition. For binary pointwise action, if
\((f(x),g(x))\in S_1\) for every \(x\), apply it at
\(u(x')\) and then apply \(v\) to both values, obtaining
\((vfu(x'),vgu(x'))\in S_2\). The negative relation does
not enter this condition; ignoring it is compatible with its
required antitone variance.

Putting whole predicates in every positive slot gives whole
constructor predicates: for functions any value lies in the
whole target, and the other constructors exhaust their carriers.
Putting order in positive binary slots gives the exact
constructor order: lifted order has bottom below everything
and old order on old values, products are coordinatewise,
nonbottom sum values compare only within their branch,
smash compares nonbottom pairs coordinatewise, and function
order is pointwise. Thus the mixed unary actions are logical
and binary actions extensional, as claimed.
For composition, apply compatibility of \(G_1\) to the
swapped relation/morphism inputs and of \(G_2\) to the
unswapped ones, then compatibility of \(F\) to those two
outputs. This is exactly the displayed \(H^{\mathrm{rel}}\)
formula. Its output is admissible because each successive
action produces an admissible relation. If each action is
logical, the positive inner output is whole, hence the outer
output is whole regardless of its negative input. If each
is extensional, the positive inner output is its constructor
order, hence the outer output is its constructor order.
This proves preservation of both properties under the specified
mixed composition, including its actual variance directions.
\end{proof}
\end{annexitem}

\begin{annexitem}{annex:p10.domain-invariance.induction-coinduction}{thm:p10-domain-invariance-induction-coinduction}{The two specializations and the countability needed for basis-term induction}
\begin{proof}
For logical induction use \(R=D\) as the whole unary relation
in the general comparison theorem. The first premise holds
because \(F^{\mathrm{rel}}(S,D)=F(D,D)\) by logicality.
The second is exactly the stated fold-closure condition.
Therefore \(D\subseteq S\), proving equality.
For coinduction let \(S=\sqsubseteq_D\). It is admissible:
it contains the bottom pair, and inequalities survive
coordinatewise suprema of increasing chains.
The first general premise is the simulation assumption.
The second holds because extensionality identifies
\(F^{\mathrm{rel}}(R,\sqsubseteq_D)\) with constructor order,
and fold is monotone. Hence \(R\subseteq\sqsubseteq_D\).
Conversely order is itself a simulation: unfold is monotone,
and extensionality gives the required constructor order relation.
Thus each ordered pair belongs to a simulation, while no simulation
contains an unordered pair, proving the iff assertion.

For the finite-expression formulation assume the compact basis
is supplied by the stated constructor syntax. The bilimit basis
formula says every compact basis point comes from a finite stage;
the primitive product/tag/lift and finite-step basis formulas at
that stage give it a finite expression. Structural induction on
these expressions proves the predicate on all compact points
when its specified constants and constructor cases are proved.
Now use omega-algebraicity. Enumerate the compact points below
\(x\), allowing repetitions if finitely many, and successively
choose a compact upper point below \(x\) for each preceding
choice and the next enumerated point. The directed compact
approximation set supplies these choices. The resulting increasing
omega-chain has supremum \(x\). Admissibility of the predicate
therefore gives it at \(x\). This use of countable bases is
under the existing ordinary countability/choice convention.
It does not assert the same consequence for an arbitrary
uncountable compact basis; the counterexample in the following
item exhibits the missing condition. Nor does this argument
prove unique normal forms for function-space expressions.
For the stream example, the two-element relation
\(\{(\bot,\bot),(r,q)\}\) is admissible: every increasing
chain in it has a largest member, unless the members agree,
in which case it is constant. The bottom pair unfolds to the
bottom pair and satisfies the lifting bottom clause.
The other pair unfolds to old pairs with equal letter \(a\)
and tails \((r,q)\) again in the relation. Thus it is a
simulation for the actual lifted product action. Coinduction
and the exchanged relation prove both inequalities and equality.
\end{proof}
\end{annexitem}

\begin{annexitem}{annex:p10.domain-invariance.failures}{thm:p10-domain-invariance-failures}{Concrete strictness, canonicality and category counterexamples}
\begin{proof}
For the identity functor the canonical stage sequence is
\(I,I,\ldots\), so its bilimit is \(I\). On the chain
\(A=\{0<1\}\), the identity structure is an isomorphism
and every continuous endomap is an algebra endomorphism.
Identity and the strict constant-zero map are different,
so it cannot be an initial strict algebra. Its least
endomorphism is constant zero, not identity.
For non-strict initiality, each map \(I\to A\) selecting
zero or one is continuous and commutes with the two identity
structures. This defeats uniqueness while the strict choice
is forced to zero.
For mixed \(F(X,Y)=Y\), take \(D=A=I\), \(B=\{0<1\}\),
\(f=1_I\), \(g=1_B\). Both structure maps are strict.
The free comparison equations permit the unique \(h:I\to I\)
and either constant \(k:I\to B\) when strictness of \(k\)
is removed. Thus backward strictness is a substantive premise.

In \(\mathcal P(U)\) compact elements are finite subsets.
All are countable. The countable subsets include empty and
are closed under unions of increasing omega-chains, by the
stated countability convention, but omit uncountable \(U\).
This is an admissible predicate containing every compact
point, not the whole domain. It defeats the unqualified
uncountable-basis inference, not the relation-action induction
theorem whose fold-closure premise is stronger.

For the set functor \(1+X\), give \(\mathbb N\) the algebra
\(\operatorname{inl}(*)\mapsto0\),
\(\operatorname{inr}(n)\mapsto n+1\).
Into any algebra \((A,a)\), its unique homomorphism is
recursively \(h(0)=a(\operatorname{inl}(*))\),
\(h(n+1)=a(\operatorname{inr}(h(n)))\). Existence is ordinary
natural-number recursion and uniqueness follows by induction.
For the final coalgebra put \(B=\mathbb N\cup\{\infty\}\)
and define its structure by
\(0\mapsto\operatorname{inl}(*)\),
\(n+1\mapsto\operatorname{inr}(n)\),
\(\infty\mapsto\operatorname{inr}(\infty)\).
For a coalgebra \(c:A\to1+A\), each point has a uniquely
determined transition path. Map it to \(n\) if it reaches
the left injection after exactly \(n\) transitions, and to
\(\infty\) if it never does. These disjoint cases define
an actual function, which satisfies the coalgebra equation.
Any homomorphism is forced to the same finite value in the
terminating case by induction on the transitions. In the
nonterminating case a finite image would force termination
after that many transitions, so its image must be infinity.
This proves finality. Its inverse structure is an algebra
on \(B\) with successor fixing infinity. Any algebra map
from the initial \(\mathbb N\) fixes each finite numeral
by induction and never hits infinity. Hence these algebra
structures are not isomorphic, even though their carriers
are both countable. This is an actual category boundary.
\end{proof}
\end{annexitem}
